\documentclass[11pt,letterpaper]{article}

\usepackage[T1]{fontenc}
\usepackage[margin=1in]{geometry}
\usepackage{parskip}
\usepackage{booktabs}
\usepackage{enumitem}
\usepackage{microtype}
\usepackage[hidelinks]{hyperref}
\usepackage{minted}

\renewcommand{\texttt}[1]{\begingroup\small\ttfamily\hyphenchar\font=45
  \spaceskip=.4em plus .2em minus .15em #1\endgroup}
\let\oldtextunderscore\_
\renewcommand{\_}{\oldtextunderscore\allowbreak}

\setminted[c]{fontsize=\small,breaklines,frame=leftline,framesep=2mm,linenos,numbersep=4pt}

\setlength{\parskip}{6pt}

\begin{document}

\begin{center}
  {\LARGE\bfseries Large Binary Filter Configuration}\\[4pt]
  {\large Filter Plugin Developer Preview --- RFC-HDFG-2026-*}\\[6pt]
  {\small M.\ Scot Breitenfeld \quad\textbullet\quad The HDF Group \quad\textbullet\quad 2026-06-10}
\end{center}

\bigskip

R.~Underwood (LibPressio) raised three issues during the
RFC-HDFG-2026-001 community review:

\begin{enumerate}[noitemsep]
  \item \textbf{Large configuration values} --- some compressors need
        multi-megabyte configuration data that cannot fit in the
        current \texttt{cd\_values} mechanism.
  \item \textbf{Non-serializable runtime settings} --- filters need
        access to process-local handles (\texttt{MPI\_Comm},
        \texttt{cudaStream\_t}, per-read quality parameters) at I/O
        time.
  \item \textbf{Chunk coordinates} --- filters need to know which
        chunk they are operating on for spatially-aware compression or
        GPU tile routing.
\end{enumerate}

\begin{center}
\small
\begin{tabular}{@{}cll@{}}
\toprule
\textbf{Issue} & \textbf{Addressed in} & \textbf{Status} \\
\midrule
1 & RFC-HDFG-2026-*  & Design in progress \\
2 & RFC-HDFG-2026-001                 & In progress \\
3 & RFC-HDFG-2026-001                 & In progress \\
\bottomrule
\end{tabular}
\end{center}

%%------------------------------------------------------------------
\section*{The Problem Being Solved}

HDF5 filter parameters are stored as an array of unsigned 32-bit integers
(\texttt{cd\_values}), capping configuration data at a few kilobytes in
practice.

RFC-HDFG-2026-001 introduced \texttt{H5Pappend\_filter} with a
human-readable key--value string; the filter's \texttt{set\_config} callback
translates that string into \texttt{cd\_values} stored on disk as before.
This removes the serialisation burden on the filter author but does not raise
the on-disk size limit.

RFC-HDFG-2026-* removes the size limit entirely for binary payloads.
The bytes are stored in the file alongside the dataset; the library retrieves
them automatically on open and hands them back to the filter.

%%------------------------------------------------------------------
\section*{Proposed Creation API}

\begin{minted}{c}
herr_t H5Pappend_filter_blob(hid_t        dcpl,
                               H5Z_filter_t filter_id,
                               unsigned int flags,
                               const void  *buf,
                               size_t       size);
\end{minted}

Call this when the filter configuration is a raw byte buffer of any size.
The bytes are stored in the file alongside the dataset (via the heap,
similar to variable-length data) and recovered automatically on open.
The performance impact is a one-time read at \texttt{H5Dopen}.

\paragraph{SZ4 JIT example:}
\begin{minted}{c}
void  *jit_buf;
size_t jit_size;
sz4_preprocess_jit_module(source_file, &jit_buf, &jit_size);

H5Pappend_filter_blob(dcpl, SZ4_FILTER_ID, H5Z_FLAG_MANDATORY,
                        jit_buf, jit_size);
free(jit_buf);   /* library holds its own copy */

H5Dcreate2(file, "mydata", dtype, space,
            H5P_DEFAULT, dcpl, H5P_DEFAULT);
\end{minted}

%%------------------------------------------------------------------
\section*{Use Cases}

\paragraph{Use Case A --- SZ4 JIT compiler (primary).}
SZ4 needs to store multi-megabyte pre-processed source code with the dataset.
This is a pure size limitation with no viable workaround under the current API.
\texttt{H5Pappend\_filter\_blob} solves it directly; the bytes travel with
the dataset and are recovered on open.

\paragraph{Use Case B --- ROIBIN-SZ spatial mask (secondary).}
ROIBIN-SZ stores a binary spatial mask indicating which values use lossless
compression.  The preferred design is a reference to a user-visible HDF5
dataset containing the mask.  Storing the mask bytes directly (the opaque
approach) makes the file self-contained but the bytes are not
user-inspectable.  Solving the reference case requires a file handle inside
the filter at open time; options are under discussion in RFC-HDFG-2026-*.

%%------------------------------------------------------------------
\section*{Issues 2 and 3 --- Addressed in RFC-HDFG-2026-001}

Issues 2 and 3 both involve information that is available at I/O time
and must reach the filter callback.  The cleanest solution for both is
to extend the filter function signature used by \texttt{H5Z\_class3\_t}
rather than add a new public API function or a separate callback.

\paragraph{Issue 2 --- Non-serializable runtime settings.}
Filters such as LibPressio need \texttt{MPI\_Comm}, \texttt{cudaStream\_t},
or per-read quality parameters that cannot be stored in \texttt{cd\_values}
because they are process-local and change between calls.

The DXPL is passed directly as a parameter to the \texttt{H5Z\_class3\_t}
filter callback.  The caller stores handles in the DXPL before I/O using
\texttt{H5Pinsert2}/\texttt{H5Pregister2} and \texttt{H5Pset}; the filter
retrieves them with \texttt{H5Pget}.  No new public API function is
required, and no file-format change is required.

\paragraph{Issue 3 --- Chunk coordinates.}
Filters need to know which chunk they are compressing or decompressing
(e.g., for spatially-varying compression or GPU tile routing).

The chunk's scaled coordinates are added as parameters to the same
extended filter callback, alongside the DXPL.  A separate callback
would introduce call-ordering hazards under multithreading; adding the
coordinates to the existing callback avoids this entirely.
No file-format change is required.


\end{document}
